\subsection{Tensor products of flat modules}

Let A, B be two rings, E a right A-module and F an (A, B)-bimodule (Algebra, Chapter II, § 1, no. 14). Recall (Algebra, Chapter II, § 3, no. 4) that \(E \otimes_{A} F\) has a canonical right B-module structure, for which

\[
(x \otimes y) b = x \otimes (y b) \quad \text {   for   } \quad x \in \mathrm{E}, y \in \mathrm{F}, b \in \mathrm{B}.
\]

PROPOSITION 8. Let A, B be two rings, E a right A-module and F an (A, B)-bimodule. Suppose that E is flat and that F is flat as a B-module. Then the B-module \(E \otimes_{A} F\) is flat.

Let G be a left B-module and \(G'\) a submodule of G. Since F is flat as a right B-module, the canonical homomorphism \(F \otimes_{B} G' \to F \otimes_{B} G\) is injective. Since E is flat, the canonical homomorphism

\[
\mathrm{E} \otimes_ {\mathbf {A}} (\mathrm{F} \otimes_ {\mathbf {B}} \mathrm{G} ^ {\prime}) \rightarrow \mathrm{E} \otimes_ {\mathbf {A}} (\mathrm{F} \otimes_ {\mathbf {B}} \mathrm{G})
\]

is injective. As \(E \otimes_{A} (F \otimes_{B} G')\) and \(E \otimes_{A} (F \otimes_{B} G)\) are canonically identified with \((E \otimes_{A} F) \otimes_{B} G'\) and \((E \otimes_{A} F) \otimes_{B} G\) respectively (Algebra, Chapter II, §3, no. 8, Proposition 8), the canonical homomorphism

\[
(\mathbf {E} \otimes_ {\mathbf {A}} \mathbf {F}) \otimes_ {\mathbf {B}} \mathbf {G} ^ {\prime} \rightarrow (\mathbf {E} \otimes_ {\mathbf {A}} \mathbf {F}) \otimes_ {\mathbf {B}} \mathbf {G}
\]

is injective, which proves that \(E \otimes_{A} F\) is a flat B-module.

COROLLARY 1. Let C be a commutative ring, E, F two flat C-modules. Then \(E \otimes_{C} F\) is a flat C-module.

F is a (C, C)-bimodule and it is sufficient to apply Proposition 8 with B = A = C.

COROLLARY 2. Let \(\rho\) be a homomorphism of a ring A to a ring B. If E is a flat right A-module, the right B-module \(\rho^{*}(\mathrm{E}) = \mathrm{E}_{(\mathrm{B})}\) obtained by extending the ring of scalars to B (Algebra, Chapter II, § 5, no. 1) is flat.

By definition \(E_{(B)} = E \otimes_{A} B\) , where B is considered as an (A, B)-bimodule by means of \(\rho\) . As the right B-module \(B_{d}\) is flat, it is sufficient to apply Proposition 8.

COROLLARY 3. Let R, S be two rings and \(\phi: \mathbf{R} \to \mathbf{S}\) a ring homomorphism. If M is a flat right S-module and \(\phi_*(\mathbf{S}_d)\) is a flat right R-module, then \(\phi_*(\mathbf{M})\) is a flat right R-module.

Recall that \(\phi_{*}(\mathbf{M})\) is the right R-module defined by \(x.r = x.\phi (r)\) for all \(x\in \mathbf{M}\) and all \(r\in \mathbb{R}\) (Algebra, Chapter II, §1, no. 13). Then apply Proposition 8 with \(\mathrm{A} = \mathrm{S}\) , \(\mathrm{B} = \mathrm{R}\) , \(\mathrm{E} = \mathrm{M}\) and \(\mathrm{F} = \mathrm{S}\) , S having the structure of a (S, R)-bimodule defined by \(\phi\) ; the right R-module M \(\otimes_{\mathbf{S}}\) S is then precisely \(\phi_{*}(\mathbf{M})\) .

PROPOSITION 9. Let \((\mathbf{A}_{\alpha}, f_{\beta \alpha})\) be a direct system of rings, \(\mathbf{A} = \varinjlim \mathbf{A}_{\alpha}\) its direct limit, \((\mathbf{E}_{\alpha}, f_{\beta \alpha})\) a direct system of right \(\mathbf{A}_{\alpha}\) -modules with the same indexing set and \(\mathbf{E} = \varinjlim \mathbf{E}_{\alpha}\) its direct limit, which is a right A-module (Algebra, Chapter II, § 6, no. 2). If each of the \(\mathbf{E}_{\alpha}\) is a flat \(\mathbf{A}_{\alpha}\) -module, \(\mathbf{E}\) is a flat A-module.

Let \(E_{\alpha}^{\prime}=E_{\alpha}\otimes_{A_{\alpha}}A\) , where A is considered as a left \(A_{\alpha}\) -module by means of the canonical homomorphism \(A_{\alpha}\to A\) ; we know that the right A-module E is canonically isomorphic to \(\varinjlim E_{\alpha}^{\prime}\) (loc. cit., Corollary 2 to Proposition 7). It follows from Corollary 2 to Proposition 8 that \(E_{\alpha}^{\prime}\) is a flat right A-module for all \(\alpha\) , hence E is a flat A-module by Proposition 2 of no. 3.

